% Part: incompleteness
% Chapter: arithmetization-syntax
% Section: introduction

\documentclass[../../../include/open-logic-section]{subfiles}

\begin{document}

\olfileid{inc}{art}{int}
\olsection{مقدمه}

برای پیوندزدن محاسبه‌پذیری و منطق، به راهی نیاز داریم تا دربارهٔ اشیای
منطق (نمادها، ترم‌ها، !!{formula}s و !!{derivation}s)، عملیات روی آن‌ها،
و ویژگی‌ها و روابطشان به شیوه‌ای مناسبِ پردازش محاسباتی سخن بگوییم.
می‌توانیم این کار را مستقیماً با درنظرگرفتن توابع و روابط محاسبه‌پذیر
روی نمادها، دنباله‌های نمادها و دیگر اشیای ساخته‌شده از آن‌ها انجام
دهیم. چون اشیای نحو منطقی همگی متناهی‌اند و از مجموعه‌های
!!a{enumerable} نمادها ساخته می‌شوند، این کار برای بعضی مدل‌های محاسبه
ممکن است. اما مدل‌های دیگری، مانند توابع بازگشتی، به اعداد و روابط و
توابع روی آن‌ها محدودند. افزون بر این، در نهایت می‌خواهیم بتوانیم
درون نظریه‌های معینی، به‌ویژه نظریه‌های صورت‌بندی‌شده در زبان حساب،
با نحو سروکار داشته باشیم. در این موارد باید نحو را \emph{حسابی
کنیم}؛ یعنی اشیای نحوی، عملیات روی آن‌ها و روابطشان را به‌ترتیب با
اعداد، توابع حسابی و روابط حسابی بازنمایی کنیم. این ایده که به
لایبنیتس بازمی‌گردد، نسبت‌دادن اعداد به اشیای نحوی است.

نسبت‌دادن اعداد به نمادها به‌عنوان «رمز»شان نسبتاً سرراست است. بعضی
نمادها اندکی دشواری ایجاد می‌کنند؛ مثلاً بی‌نهایت !!{variable}s و حتی
برای هر آریتیِ~$n$ بی‌نهایت !!{function}s وجود دارد. اما البته می‌توان
اعداد را به‌طور منظم چنان به نمادها نسبت داد که مثلاً $\Obj v_2$ و
$\Obj v_3$ رمزهای متفاوت داشته باشند. دنباله‌های نمادها (مانند ترم‌ها
و !!{formula}s) دشوارترند. ولی اگر بتوانیم با دنباله‌های اعداد به‌طور
کاملاً حسابی کار کنیم (مثلاً با رمزگذاری دنباله‌ها به توان‌های اعداد
اول)، می‌توانیم رمزگذاری نمادهای منفرد را به دنباله‌های نمادها و سپس
به دنباله‌ها یا آرایش‌های دیگری از !!{formula}s، مانند !!{derivation}s،
گسترش دهیم. این رمزگذاری گسترش‌یافته «شماره‌گذاری گودل» نام دارد. به
هر ترم، !!{formula} و !!{derivation} یک عدد گودل نسبت داده می‌شود.

با رمزگذاری دنباله‌های نمادها به‌عنوان دنباله‌های رمزهایشان و انتخاب
دستگاهی برای رمزگذاری دنباله‌ها که با توابع محاسبه‌پذیر پردازش‌پذیر
باشد، می‌توانیم اعداد گودل را نیز با توابع محاسبه‌پذیر پردازش کنیم.
در عمل، همهٔ توابع مربوط بازگشتی اولیه خواهند بود. برای نمونه، هم
محاسبهٔ طول یک دنباله و هم محاسبهٔ عضو~$i$اُم آن از روی رمز دنباله
بازگشتی اولیه‌اند. اگر عدد رمزکنندهٔ دنباله مثلاً عدد گودلِ
!!a{formula}~$!A$ باشد، بی‌درنگ می‌بینیم طول !!a{formula} و (رمزِ) نماد
$i$اُم در !!a{formula} را نیز می‌توان از عدد گودل~$!A$ محاسبه کرد. اثبات اینکه،
مثلاً، ویژگیِ عدد گودلِ یک ترم خوش‌ساخت یا !!{derivation} درست بودن
بازگشتی اولیه است اندکی دشوارتر است. بااین‌حال ممکن است، زیرا
دنباله‌های مورد نظر (ترم‌ها، !!{formula}s و !!{derivation}s) به‌طور
استقرایی تعریف شده‌اند.

به‌عنوان نمونه، عمل جانشانی را در نظر بگیرید. اگر $!A$ فرمول، $x$
متغیر و~$t$ ترم باشد، $\Subst{!A}{t}{x}$ حاصل جایگزین‌کردن هر رخداد
آزاد~$x$ در~$!A$ با~$t$ است. اکنون فرض کنید اعداد گودل~$k$، $l$ و~$m$
را به‌ترتیب به~$!A$، $x$ و~$t$ نسبت داده‌ایم. همان طرح به
$\Subst{!A}{t}{x}$ نیز عدد گودلی، مثلاً~$n$، نسبت می‌دهد. این نگاشت
از~$k$، $l$ و~$m$ به~$n$ همتای حسابیِ عمل جانشانی است. وقتی عمل
جانشانی $!A$، $x$ و~$t$ را به~$\Subst{!A}{t}{x}$ می‌فرستد، تابع
حسابی‌شدهٔ جانشانی اعداد گودل~$k$، $l$ و~$m$ را به عدد گودل~$n$
می‌فرستد. خواهیم دید این تابع بازگشتی اولیه است.

حسابی‌سازی نحو فقط از نظر انتزاعی جالب نیست، هرچند در آغاز این بینش
که زبان‌هایی مانند زبان حساب، که سازوکاری برای «سخن‌گفتن دربارهٔ»
زبان‌ها ندارند، بااین‌همه می‌توانند ویژگی‌های پیچیدهٔ عبارت‌ها را
صوری کنند غیرپیش‌پاافتاده بود. از آنجا فقط گامی کوچک است تا بپرسیم
نظریه‌ای در این زبان، مانند حساب پئانو، چه چیزهایی را دربارهٔ زبان
خود \emph{اثبات} می‌کند (از جمله اینکه آیا !!{sentence}s اثبات‌پذیر یا
صادق‌اند). این ما را به قضایای تحدیدی مشهور گودل (دربارهٔ
اثبات‌ناپذیری) و تارسکی (دربارهٔ تعریف‌ناپذیری صدق) می‌رساند. اما
ترفند حسابی‌سازی نحو برای اثبات بعضی نتایج مهم نظریهٔ محاسبه‌پذیری،
مثلاً دربارهٔ توان محاسباتی نظریه‌ها یا رابطهٔ مدل‌های گوناگون
محاسبه‌پذیری، نیز مهم است. حسابی‌سازی نحو الگویی برای حسابی‌سازی دیگر
اشیا و ویژگی‌هاست. برای نمونه، به همین ترتیب می‌توان پیکربندی‌ها و
محاسبات (مثلاً ماشین‌های تورینگ) را حسابی کرد. این امکان می‌دهد
محاسبات یک مدل (مثلاً ماشین‌های تورینگ) را در مدلی دیگر (مثلاً توابع
بازگشتی) شبیه‌سازی کنیم.

\end{document}
